\cvsection{Expériences professionnelles}

\begin{cventries}
  \cventry
    {Stagiaire ingénieur de recherche en IA}
    {INRIA, équipe Mnémosyne}
    {Bordeaux, France}
    {Mars 2026 -- Septembre 2026}
    {
      \begin{cvitems}
        \item {Conception et développement en autonomie, avec le suivi de mes encadrants, d’une application aidant les chercheurs à trouver des articles scientifiques pertinents.}
        \item {Développement d’un outil Python de collecte automatisée de métadonnées scientifiques à partir de sources publiques.}
        \item {Mise en œuvre d’un pipeline RAG avec découpage des documents, embeddings, recherche sémantique et stockage vectoriel dans ChromaDB.}
        \item {Conception d’un agent IA multi-étapes, de ses outils et de ses prompts pour rechercher, sélectionner et restituer les informations utiles au cas d’usage.}
        \item {Développement de services Python/FastAPI, d’API REST documentées avec OpenAPI et de la persistance MongoDB.}
        \item {Développement d’une interface React/Next.js permettant aux chercheurs d’interroger l’application et d’exploiter ses résultats.}
        \item {Comparaison de Qwen, Gemma et GLM avec plusieurs configurations de quantification sur une infrastructure HPC avec Slurm.}
        \item {Définition d’un protocole expérimental et analyse de la latence, du TTFT, de la qualité des réponses et de l’effet de la taille du contexte.}
        \item {Contribution à un article scientifique documentant la méthode, les résultats et les limites observées.}
        \item {Conteneurisation de composants avec Docker et mise en place d’un pipeline CI/CD sur GitHub.}
      \end{cvitems}
    }

  \cventry
    {Stagiaire de recherche en apprentissage par renforcement}
    {Université Karunya}
    {Coimbatore, Tamil Nadu, Inde}
    {Juin 2025 -- Septembre 2025}
    {
      \begin{cvitems}
        \item {Développement d’un modèle d’apprentissage par renforcement et de son environnement d’entraînement pour la navigation autonome d’un robot hospitalier.}
        \item {Déploiement et optimisation du logiciel et du modèle sur un Raspberry Pi embarqué.}
      \end{cvitems}
    }
\end{cventries}
